\section{Nesso-1 architecture: a high-level view}

Nesso-1 asks whether binding affinity can be estimated from an approximate map
of the protein--ligand interface without first generating every atomic
coordinate. It is therefore best understood as a \emph{coarse-grained
cofolding affinity model}, rather than as a full-atom structure predictor
\citep{shenoy2026nesso1}.

\begin{center}
\fbox{\parbox{0.92\textwidth}{\centering
Protein sequence + ligand SMILES
$\longrightarrow$ protein and ligand embeddings
$\longrightarrow$ pairwise distance representation
$\longrightarrow$ likely binding-region crop
$\longrightarrow$ binder probability + continuous affinity score}}
\end{center}

\subsection{Inputs and molecular tokens}

The inputs are a protein amino-acid sequence and a ligand SMILES string. No
experimental complex, user-specified pocket, or multiple-sequence alignment is
required. The protein receives contextual embeddings from the pretrained
ESM-2 protein language model. Each protein residue is represented
geometrically by one center, its C$\beta$ atom (C$\alpha$ for glycine), while
each ligand heavy atom remains an individual token. Thus, the protein is
strongly coarse-grained but the ligand retains its heavy-atom graph.

\subsection{A learned map of pairwise geometry}

The central object is a pair representation $z$. For every pair of tokens,
$z_{ij}$ stores a learned description of their relationship. Protein--ligand
entries answer questions such as ``is residue 42 likely to be close to ligand
atom 5?'' A streamlined PairFormer updates this matrix using triangle
attention and triangle multiplication, operations that allow relationships
among three tokens to constrain one another.

The trunk predicts a categorical distance distribution,
\[
  P(d_{ij}\in\text{distance bin }k),
\]
rather than one exact coordinate set. From this distribution, Nesso-1 obtains
an expected distance and an entropy that measures uncertainty in the predicted
protein--ligand geometry.

\subsection{Finding the relevant protein region}

The first trunk pass considers the full protein. Subsequent recycling steps
retain protein tokens predicted to lie within approximately 22~\AA{} of the
ligand. The affinity module uses a tighter crop: residues predicted to lie
within approximately 15~\AA{} of any ligand heavy atom. In this sense, Nesso-1
infers a likely interaction region internally rather than requiring the user to
provide a pocket.

\subsection{Affinity outputs}

A second, smaller PairFormer---the affinity module---combines the protein
language-model embeddings, ligand features, pair representation $z$, and its
predicted distance information. It produces two distinct outputs:

\begin{enumerate}
  \item a binder likelihood for hit-identification classification; and
  \item a continuous affinity or potency score for compound ranking.
\end{enumerate}

The continuous output is not a calculated $\Delta G_{\mathrm{bind}}$. During
training, $K_i$, $K_d$, IC$_{50}$, and EC$_{50}$ labels are treated
interchangeably to accommodate incomplete and heterogeneous public assay
metadata. It should therefore be interpreted as an assay-trained
affinity-like score, with particular value for ranking compounds measured in
comparable settings.

\subsection{The essential contrast with Boltz-2}

\begin{center}
\begin{tabular}{p{0.21\textwidth}p{0.34\textwidth}p{0.34\textwidth}}
\toprule
 & \textbf{Boltz-2} & \textbf{Nesso-1} \\
\midrule
Structural decoder & Full-atom diffusion producing explicit coordinates & No
heavy-atom diffusion; predicts coarse pair distances \\
Protein context & MSA can be used & ESM-2 embeddings; no MSA \\
Primary structural output & Full biomolecular complex & Coarse interface
representation; approximate poses can be reconstructed \\
Design priority & Detailed structure plus affinity & Rapid affinity screening \\
\bottomrule
\end{tabular}
\end{center}

Removing full-atom diffusion is the main source of Nesso-1's speed. The cost is
that it cannot resolve side-chain conformations, hydrogen-bond geometry,
steric clashes, or water-mediated interactions with atomistic precision. Its
architectural hypothesis is that coarse interface geometry may be sufficient
for rapid affinity ranking, even though full physical detail remains important
for mechanistic interpretation and late-stage optimization.
